TODO Something which gets us onto the package track.

Packages come in many different forms, and what they are mostly depends on the
system that they serve. For example, binary package managers may ship
applications which are then installed on the user’s system. An example of
such as package manager could be the iOS App Store. In general, a package
manager primarily makes several things easier:

\begin{enumerate}
    \item Publishing of Packages
    \item Find packages to install
    \item Automatically install dependencies
    \item Handle available updates
\end{enumerate}

These same principles have also been applied to the application level.
Application-level packages are usually software components, which are intended
to be run locally as a part of your application. Package managers also
typically attempt to establish some sort of language/framework convention for
developing packages. This is done such that packages may be imported into a
software project in a consistent manner.

The goal of this thesis is to develop a package manager for the
Jolie\autocite{montesi2010jolie,JOLA} ecosystem. Jolie is a service-oriented
programming language.
% Do we really need this?
As such, we aim not only at covering the local packages (i.e., services running
        locally), but also external services. From the perspective of a local
Jolie service, an external service is any service, which runs on an external
virtual machine.

Before being able to achieve the goal of creating a package manager, a few
hurdles must first be overcome. In this thesis, we will present a new natively
supported module and configuration system for Jolie.

% Should also cover the somewhat blurry line there is between package manager
% and build tool

% Method
% ------
% This section focuses a bit too much on looking at other products. We should
% still mention a few up here. But we never actually looked straight at other
% products.
%
% There is also no focus on the configuration system which is developed.

\pagebreak

\section{Structure of this Thesis}

\textbf{Chapter 1}. This chapter. Introduces the thesis, explains the structure
of it, and notations used throughout the thesis.

\textbf{Chapter 2}. Covers background material which is relevant for the entire
thesis.

\textbf{Chapter 3}. Covers the module system and the configuration system which
had been made to support the package manager efforts.

\textbf{Chapter 4}. Explains the concept of ``Jolie Packages'' which are what
the Jolie Package Manager (JPM) supports. This chapter will also cover many
of the features that are directly related to packages, which JPM implements.

\textbf{Chapter 5}. Covers architectural and implementation specific details
for JPM.

\textbf{Chapter 6}. Conclusion of this thesis.

\section{Notation Used}

\subsection*{Service Diagrams}

Simple box-and-line diagrams are used for illustrating how service are
deployed, and how they depend on eachother. Figure
\ref{fig:simple_box_and_line} shows a simple diagram with two services \txtl{A}
and \txtl{B}. In this diagram \txtl{A} depends on the services provided by
\txtl{B}.

\begin{figure}[H]

\centering
\includegraphics[width=0.5\textwidth]{introduction/simple_box_and_line.eps}

\caption{A service diagram showing two services: \txtl{A} and \txtl{B}. In this
    diagram service \txtl{A} depends on service \txtl{B}}

\label{fig:simple_box_and_line}

\end{figure}

We'll illustrate embedding (a Jolie concept, see Chapter \ref{cha:background})
by drawing a service inside another. The parent service will show its
dependency on this service through similar means.

\begin{figure}[H]

\centering
\includegraphics[width=0.5\textwidth]{introduction/embedding.eps}

\caption{A \txtl{parent} service embedding a \txtl{child} service}

\label{fig:box_and_line_embedding}

\end{figure}

\subsection*{File Listings}

When code snippets are presented, where the file name and location is relevant
then the file name and file path will be shown in a comment before the file
contents. We will commonly show multiple files in a single listing.  As an
example, if we have the directory structure shown in Listing
\ref{lst:file_structure_example}.

\begin{listing}[H]
\begin{minted}{text}
/
|-- a
|   `-- 1
|-- b
|   `-- 2
`-- c
    `-- 3

3 directories, 3 files
\end{minted}
\caption{We will illustrate a directory structure in this manner}
\label{lst:file_structure_example}
\end{listing}

And the file contents of file \txtl{/a/1} is ``A'', \txtl{/b/2} contains ``B'',
    and \txtl{/c/3} contains ``C'', then we would show the contents of these
    files as in Listing \ref{lst:multiple_files_example}.

\begin{listing}[H]
\begin{minted}{text}
/* /a/1 */
A

/* /b/2 */
B

/* /c/3 */
C
\end{minted}

\caption{Multiple files shown in a single listing where file path and name is
    relevant.}

\label{lst:multiple_files_example}
\end{listing}

\subsection*{Command Line Illustrations}

In certain cases, examples are made by showing the output of a command line
session. The reader is expected to be familiar with a basic Unix-like shell.
When the working directory is relevant, then it will be shown as part of the
prompt. An example of this is shown in Listing \ref{lst:shell_example}.  In
this example two commands are invoked (\txtl{ls} and \txtl{cat hello-world.ol}.
The \verb!#! symbols mark the beginning of a command. Lines not
containing this symbol should be considered as part of the output from
the previous command. Before that the absolute path to the current
working directory is shown, in these cases they were both
\txtl{/this/is/the/path}.

\begin{listing}[H]
\begin{minted}{text}
dan@host:/this/is/the/path # ls
hello-world.ol

dan@host:/this/is/the/path # cat hello-world.ol
include "console.iol"
main {
    println@Console("Hello, World!")()
}
\end{minted}
\caption{Illustration of two commands being invoked in a shell}
\label{lst:shell_example}
\end{listing}

\section{Implementation Status}

